\documentclass[11pt]{article}




\usepackage[sfdefault]{FiraSans} %% option 'sfdefault' activates Fira Sans as the default text font
\usepackage[T1]{fontenc}
\renewcommand*\oldstylenums[1]{{\firaoldstyle #1}}

\usepackage{natbib}
\usepackage[french,english]{babel}
\usepackage{numprint}
\usepackage{multirow}
\usepackage{rotating}
\usepackage{fancyhdr}
\usepackage{booktabs}
\usepackage{multicol}
\usepackage{hyperref}\hypersetup{colorlinks=true}

\usepackage{amsmath,amssymb,amsfonts,textcomp}
\usepackage{color}
\usepackage{calc}
 \setlength{\tabcolsep}{8pt}
\usepackage{setspace}
\onehalfspacing
\usepackage{longtable}
\usepackage{graphicx}
\usepackage[margin=1in]{geometry}
\setlength{\parindent}{0pt}
\usepackage[bottom]{footmisc}
\pagestyle{fancy}
\usepackage{titlesec}
\usepackage{lipsum}
\usepackage{cancel}
\usepackage{multicol}

\usepackage{amsmath,amssymb}
\usepackage{lmodern}
\usepackage{iftex}
\ifPDFTeX
  \usepackage[T1]{fontenc}
  \usepackage[utf8]{inputenc}
  \usepackage{textcomp} % provide euro and other symbols
\else % if luatex or xetex
  \usepackage{unicode-math}
  \defaultfontfeatures{Scale=MatchLowercase}
  \defaultfontfeatures[\rmfamily]{Ligatures=TeX,Scale=1}
\fi

\titleformat{\section}
  {\normalfont\Large\scshape\bfseries}{\thesection}{1em}{}
  \titlespacing{\section}{0pt}{10pt}{0pt}
\titleformat{\subsection}
  {\normalfont\bfseries}{\thesection}{1em}{}
  \titlespacing{\subsection}{0pt}{6pt}{0pt}

\lhead{ECON3500 - Econometrics and Applications}

\setlength\parskip{0.10in}
\begin{document}
\thispagestyle{plain}
\singlespacing


Version: \today \hfill \\
ECON3500: Econometrics and Applications
%\vspace{1.5cm}
\begin{center}
\Large{\textbf{Problem Set 4}}\\
\end{center}
\bigskip


\hypertarget{welcome}{%
\section*{Welcome}\label{welcome}}

Chapters 8 and 9 problems! Enjoy!

See the exercises below, or you can \href{https://econ3500s26.netlify.app/assignment/04-ps.pdf}{download them as
a pdf}.

\hypertarget{what-do-i-submit}{%
\section*{What do I submit?}\label{what-do-i-submit}}

\begin{itemize}
\item
  Your written up answers to exercise questions. If you work on a piece
  of paper, please scan using some sort of phone software (like
  Microsoft Lens or Adobe Scan) rather than just taking a picture.
\item
  A do-file that runs your Stata analysis (for question 8).
\item
  A log file that includes the output from running your do-file (for
  question 8).
\end{itemize}

\hypertarget{exercises}{%
\section*{Exercises}\label{exercises}}

\begin{enumerate}
\def\labelenumi{\arabic{enumi}.}

\item
  The following equation describes the median housing price in a
  community in terms of amount of pollution (\(nox\) for nitrous oxide)
  and the average number of rooms in houses in the community
  (\(rooms\)):


\(log(price) = \beta_0 + \beta_1log(nox) + \beta_2rooms + u\)

\begin{enumerate}
\def\labelenumi{\alph{enumi}.}

\item
  What are the probable signs of \(\beta_1\) and \(\beta_2\)? What is
  the interpretation of \(\beta_1\)? Explain.
\item
  Why might \(nox\) {[}or more precisely, \(log(nox)\){]} and \(rooms\)
  be negatively correlated? If this is the case, does the simple
  regression of \(log(price)\) on \(log(nox)\) produce an upward or a
  downward biased estimator of \(\beta_1\)?
\item
  Using data, the following equations were estimated:
\end{enumerate}

\(\widehat{log(price)} = 11.71 - 1.043 log(nox)\), \(n = 506\),
\(R^2 = 0.264\)

\(\widehat{log(price)} = 9.23 - 0.718 log(nox) + 0.306 rooms\),
\(n = 506\), \(R^2 = 0.514\)

Is the relationship between the simple and multiple regression estimates
of the elasticity of \(price\) with respect to \(nox\) what you would
have predicted, given your answer in part (b)? Does this mean that
0.718 is definitely closer to the true elasticity than 1.043?


\item   Read the box \emph{``The Return to Education and the Gender Gap''} in
  Section 8.3 of your textbook (Stock \& Watson).


\begin{enumerate}
\def\labelenumi{\alph{enumi}.}

\item
  Consider a man with 16 years of education and 2 years of experience.
  Use the results from column (4) of Table 8.1 and the method in Key
  Concept 8.1 to estimate the expected change in the logarithm of
  average hourly earnings (AHE) associated with an additional year of
  experience.
\item
  Explain why your answer to (a) does not depend on the region he is
  from.
\item
  Repeat (a), assuming 10 years of experience.
\end{enumerate}


\item
  To answer this question, refer to \emph{Table 8.3: Nonlinear
  Regression Model of Test Scores} in your textbook:


\begin{enumerate}
\def\labelenumi{\alph{enumi}.}

\item
  A researcher suspects that the effect of \% Eligible for subsidized
  lunch has a nonlinear effect on test scores. In particular, he
  conjectures that increases in this variable from 10\% to 20\% have
  little effect on test scores but that changes from 50\% to 60\% have a
  much larger effect. i. Describe a nonlinear specification that can be
  used to model this form of nonlinearity. ii. How would you test
  whether the researcher's conjecture was better than the linear
  specification in column (7) of Table 8.3?
\item
  A researcher suspects that the effect of income on test scores is
  different in districts with small classes than in districts with large
  classes. i. Describe a nonlinear specification that can be used to
  model this form of nonlinearity.
\end{enumerate}



\item
  Labor economists studying the determinants of women's earnings
  discovered a puzzling empirical result. Using randomly selected
  employed women, they regressed earnings on the women's number of
  children and a set of control variables (age, education, occupation,
  and so forth). They found that women with more children had higher
  wages, controlling for these other factors. Explain how sample
  selection might be the cause of this result. (Hint: Notice that women
  who do not work outside the home are missing from the sample.) {[}This
  empirical puzzle motivated James Heckman's research on sample
  selection that led to his 2000 Nobel Prize in Economics. See Heckman
  (1974){]}


\item
  This question uses directed acyclic graphs (DAGs), which we will cover in class. You may also find it helpful to read Huntington-Klein, \emph{The Effect}, \href{https://theeffectbook.net/ch-CausalPaths.html}{Chapter 8: Causal Paths and Closing Back Doors}, especially Sections 8.3--8.5.

  Consider the relationship between a woman's number of children and her earnings from question 4.

\begin{enumerate}
\def\labelenumi{\alph{enumi}.}

\item
  Draw a DAG that includes the following variables: Earnings, Number of Children, Decision to Work Outside the Home, and Ability/Motivation. Add arrows representing plausible causal relationships. For each arrow, write one sentence explaining why you included it.
\item
  Is ``Decision to Work Outside the Home'' a confounder, a collider, or a mediator on the path between Number of Children and Earnings? Explain.
\item
  When researchers study only employed women, they are conditioning on ``Decision to Work.'' Using your DAG, explain why this could produce a spurious positive relationship between number of children and earnings --- even if children have no direct causal effect on earnings.
\end{enumerate}


\item
  The demand for a commodity is given by
  \(Q = \beta_0 + \beta_1 P + u\), where \(Q\) denotes quantity, \(P\)
  denotes price, and \(u\) denotes factors other than price that
  determine demand. Supply for the commodity is given by
  \(Q = \gamma_0 + \gamma_1P + v\), where \(v\) denotes factors other
  than price that determine supply. Suppose \(u\) and \(v\) both have a
  mean of 0, have variances \(\sigma^2_u\) and \(\sigma^2_v\), and are
  mutually uncorrelated.

\begin{enumerate}
\def\labelenumi{\alph{enumi}.}

\item
  Solve the two simultaneous equations to show how Q and P depend on u
  and v. \emph{(Hint: In equilibrium, quantity supplied equals quantity demanded. Set the two equations equal and solve for \(P\) in terms of \(u\) and \(v\). Then substitute back to find \(Q\).)}
\item
  Derive the means of P and Q. \emph{(Hint: Use your answers from part (a) and the fact that \(E(u) = E(v) = 0\).)}
\item
  \textbf{(Optional)} Derive the variance of P, the variance of Q, and the covariance
  between Q and P.
\end{enumerate}


\item
  Revisit the box \emph{``The Return to Education and the Gender Gap''}
  in Section 8.3 of your textbook (Stock \& Watson). Discuss the internal and external validity of the
  estimated effect of education on earnings.

\item
  Use the dataset
  \href{https://www.princeton.edu/~mwatson/Stock-Watson_3u/Students/EE_Datasets/CollegeDistance.dta}{\texttt{CollegeDistance.dta}}
  (described in Empirical Exercise AEE 4.3) to answer the following questions.

\begin{enumerate}
\def\labelenumi{\alph{enumi}.}

\item
  Run a regression of \(ED\) on \(Dist\), \(Female\), \(Bytest\), \(Tuition\), \(Black\), \(Hispanic\), \(Incomehi\), \(Ownhome\), \(DadColl\), \(MomColl\), \(Cue80\), and \(Stwmfg80\). If \(Dist\) increases from 2 to 3 (that is, from 20 to 30 miles), how are years of education expected to change? If \(Dist\) increases from 6 to 7 (that is, from 60 to 70 miles), how are years of education expected to change?

\item
  Run a regression of \(\ln(ED)\) on \(Dist\), \(Female\), \(Bytest\), \(Tuition\), \(Black\), \(Hispanic\), \(Incomehi\), \(Ownhome\), \(DadColl\), \(MomColl\), \(Cue80\), and \(Stwmfg80\). If \(Dist\) increases from 2 to 3 (from 20 to 30 miles), how are years of education expected to change? If \(Dist\) increases from 6 to 7 (from 60 to 70 miles), how are years of education expected to change?

\item
  Run a regression of \(ED\) on \(Dist\), \(Dist^2\), \(Female\), \(Bytest\), \(Tuition\), \(Black\), \(Hispanic\), \(Incomehi\), \(Ownhome\), \(DadColl\), \(MomColl\), \(Cue80\), and \(Stwmfg80\). If \(Dist\) increases from 2 to 3 (from 20 to 30 miles), how are years of education expected to change? If \(Dist\) increases from 6 to 7 (from 60 to 70 miles), how are years of education expected to change?

\item
  Do you prefer the regression in (c) to the regression in (a)? Explain.

% \item
%   Consider a Hispanic female with \(Tuition = \$950\), \(Bytest = 58\), \(Incomehi = 0\), \(Ownhome = 0\), \(DadColl = 1\), \(MomColl = 1\), \(Cue80 = 7.1\), and \(Stwmfg = \$10.06\).
%   \begin{enumerate}
%   \def\labelenumii{\roman{enumii}.}
%   \item Plot the regression relation between \(Dist\) and \(ED\) from (a) and (c) for \(Dist\) in the range of 0 to 10 (from 0 to 100 miles). Describe the similarities and differences between the estimated regression functions. Would your answer change if you plotted the regression function for a white male with the same characteristics?
%   \item How does the regression function (c) behave for \(Dist > 10\)? How many observations are there with \(Dist > 10\)?
%   \end{enumerate}

\item
  Add the interaction term \(DadColl \times MomColl\) to the regression in (c). What does the coefficient on the interaction term measure?

\item
  Mary, Jane, Alexis, and Bonnie have the same values of \(Dist\), \(Bytest\), \(Tuition\), \(Female\), \(Black\), \(Hispanic\), \(Incomehi\), \(Ownhome\), \(Cue80\), and \(Stwmfg80\). Neither of Mary's parents attended college. Jane's father attended college, but her mother did not. Alexis's mother attended college, but her father did not. Both of Bonnie's parents attended college. Using the regressions from (e):
  \begin{enumerate}
  \def\labelenumii{\roman{enumii}.}
  \item What does the regression predict for the difference between Jane's and Mary's years of education?
  \item What does the regression predict for the difference between Alexis's and Mary's years of education?
  \item What does the regression predict for the difference between Bonnie's and Mary's years of education?
  \end{enumerate}

\item
  Is there any evidence that the effect of \(Dist\) on \(ED\) depends on the family's income?

\item
  After running all these regressions (and any others that you want to run), summarize the effect of \(Dist\) on years of education.
\end{enumerate}
\end{enumerate}

\end{document}
